\section{Value Learning: Deep Q-Networks (DQN)}

\subsection{Idea básica de DQN}

Deep Q-Network (DQN) es el algoritmo representativo de Value Learning. Su idea central es: usar una red neuronal profunda para \textbf{aproximar la Q-function}, y después usar la Q-function aprendida para derivar la política óptima.

\begin{figure}[H]
\centering
\includegraphics[width=0.95\textwidth]{slides-images/slide-26.jpg}
\caption{Las dos arquitecturas de DQN: a la izquierda, (state, action) $\to$ Q-value; a la derecha, state $\to$ Q-value de todas las acciones\protect\footnotemark}
\end{figure}
\footnotetext{Slides, página 26.}

DQN tiene dos diseños de arquitectura equivalentes:

\begin{enumerate}
    \item \textbf{Opción uno}: entrada $(s, a)$, salida un único valor $Q(s,a)$. Se necesita llamar a la red una vez por cada acción posible.
    \item \textbf{Opción dos (más eficiente)}: entrada $s$, salida simultánea de los valores Q de todas las acciones $Q(s, a_1), Q(s, a_2), \ldots, Q(s, a_n)$. Basta con una sola pasada hacia adelante.
\end{enumerate}

En la práctica normalmente se usa la opción dos, porque es más eficiente: con una sola llamada a la red se obtiene la evaluación de valor de todas las acciones.

\subsection{Proceso de decisión de DQN}

\begin{figure}[H]
\centering
\includegraphics[width=0.85\textwidth]{slides-images/slide-36.jpg}
\caption{Ejemplo de decisión de DQN: se elige la acción con el valor Q más alto $a_1$ (moverse hacia la izquierda)\protect\footnotemark}
\end{figure}
\footnotetext{Slides, página 36.}

El proceso de decisión es muy directo:
\begin{enumerate}
    \item Se introduce el estado actual $s$ (por ejemplo, la pantalla de un juego de Atari) en DQN
    \item La red genera el valor Q de cada acción posible
    \item Se elige y se ejecuta la acción con el valor Q más grande: $\pi(s) = \arg\max_a Q(s, a)$
\end{enumerate}

\subsection{Ejemplo de Atari Breakout: entender el valor Q}

En clase se usó el juego Atari Breakout para profundizar en el significado de la Q-function. Se consideran dos pares state-action:

\begin{figure}[H]
\centering
\includegraphics[width=0.95\textwidth]{slides-images/slide-22.jpg}
\caption{Comparación de valores Q en Atari Breakout: ¿cuál par (s, a) tiene el valor Q más alto?\protect\footnotemark}
\end{figure}
\footnotetext{Slides, página 22.}

\begin{itemize}
    \item \textbf{Escenario A}: la pelota cae en línea recta hacia el centro de la paleta, y la paleta permanece inmóvil
    \item \textbf{Escenario B}: la pelota vuela hacia un costado y la paleta se mueve a la derecha para alcanzarla
\end{itemize}

\begin{knowledgebox}{Intuición sobre el valor Q}
Intuitivamente el valor Q del escenario B podría parecer más alto (la pelota rebota hacia el costado y golpea más ladrillos), pero los experimentos reales muestran que la estrategia del escenario B (el golpe lateral), aunque puede dar una puntuación mayor en una sola jugada, deja la paleta desplazada hacia el costado y hace que la velocidad de devolución sea más lenta, lo que a largo plazo resulta desfavorable. Esto muestra que el valor Q mide el \textbf{retorno acumulado a largo plazo} y no la recompensa de un solo paso; la intuición humana puede no estimar con precisión el valor Q, y precisamente por eso se necesita una red profunda para aprenderlo.
\end{knowledgebox}

\subsection{Entrenamiento de DQN}

El núcleo del entrenamiento de DQN es definir un valor objetivo y una función de pérdida adecuados.

\textbf{Target} (valor objetivo): se construye el objetivo de entrenamiento a partir de la ecuación de Bellman:

$$\text{target} = r + \gamma \max_{a'} Q(s', a')$$

\begin{itemize}
    \item $r$: la recompensa inmediata obtenida en el paso actual
    \item $s'$: el nuevo estado al que se llega tras ejecutar la acción
    \item $\max_{a'} Q(s', a')$: el valor Q más grande entre todas las acciones posibles en el nuevo estado (suponiendo que después se sigue la política óptima)
\end{itemize}

\textbf{Q-Loss} (función de pérdida):

$$\mathcal{L} = \mathbb{E}\left[\left\| \underbrace{\left(r + \gamma \max_{a'} Q(s', a')\right)}_{\text{target}} - \underbrace{Q(s, a)}_{\text{predicted}} \right\|^2\right]$$

\begin{figure}[H]
\centering
\includegraphics[width=0.95\textwidth]{slides-images/slide-30.jpg}
\caption{Entrenamiento de DQN: la Q-Loss se define como el error cuadrático medio entre el valor Q target y el predicted\protect\footnotemark}
\end{figure}
\footnotetext{Slides, página 30.}

\begin{importantbox}{Intuición central del entrenamiento de DQN}
El entrenamiento de DQN consiste, en esencia, en hacer que el valor Q predicho por la red se acerque gradualmente al valor target dado por la ecuación de Bellman. El target está formado por «la recompensa inmediata más el retorno futuro óptimo descontado», con lo que se logra el objetivo de «usar el retorno futuro para guiar la decisión actual». El proceso de entrenamiento realiza la retropropagación y la actualización de parámetros minimizando el error cuadrático medio entre el Q predicted y el Q target.
\end{importantbox}

\subsection{Resultados de DQN en juegos de Atari}

Google DeepMind publicó en 2015 en Nature el artículo histórico de DQN, en el que se usó la misma arquitectura de red y los mismos hiperparámetros para hacer pruebas en 49 juegos de Atari.

\begin{figure}[H]
\centering
\includegraphics[width=0.95\textwidth]{slides-images/slide-32.jpg}
\caption{Arquitectura de red de DQN: las capas convolucionales procesan la pantalla del juego y las capas totalmente conectadas generan el valor Q de cada acción\protect\footnotemark}
\end{figure}
\footnotetext{Slides, página 32. Mnih y colaboradores, Nature 2015.}

La arquitectura de red incluye:
\begin{itemize}
    \item Varias capas de \textbf{Convolution} (capas convolucionales): procesan los píxeles de la pantalla del juego sin procesar
    \item Varias capas \textbf{Fully Connected} (totalmente conectadas): mapean las características al valor Q de cada acción
    \item Capa de salida: cada nodo corresponde a una acción posible (por ejemplo, cada combinación de direcciones y botones del control del juego)
\end{itemize}

\begin{figure}[H]
\centering
\includegraphics[width=0.95\textwidth]{slides-images/slide-33.jpg}
\caption{Desempeño de DQN en juegos de Atari: más del 50\% de los juegos alcanzan o superan el nivel humano\protect\footnotemark}
\end{figure}
\footnotetext{Slides, página 33. Mnih y colaboradores, Nature 2015.}

\begin{knowledgebox}{Importancia histórica de DQN}
La importancia de DQN radica en que demostró que un sistema de aprendizaje por refuerzo profundo \textbf{de extremo a extremo} puede aprender directamente, a partir de píxeles sin procesar, políticas complejas. Antes de esto, el RL solía requerir el diseño manual de características (Feature Engineering). DQN abrió una nueva era del Deep RL y fue una de las tecnologías clave que llevaron a que Google adquiriera DeepMind.
\end{knowledgebox}

\subsection{Limitaciones de Q-Learning}

\begin{figure}[H]
\centering
\includegraphics[width=0.85\textwidth]{slides-images/slide-34.jpg}
\caption{Dos grandes limitaciones de Q-Learning: no puede manejar espacios de acción continuos y solo puede aprender políticas deterministas\protect\footnotemark}
\end{figure}
\footnotetext{Slides, página 34.}

Aunque DQN logró un enorme éxito, el método Q-Learning tiene dos limitaciones importantes:

\textbf{Problema de complejidad} (Complexity):
\begin{itemize}
    \item Solo puede manejar espacios de acción \textbf{discretos y finitos}
    \item Para espacios de acción continuos (como controlar el ángulo del volante), no es posible enumerar todas las acciones posibles
\end{itemize}

\textbf{Problema de flexibilidad} (Flexibility):
\begin{itemize}
    \item La política se deriva de la Q-function de manera determinista mediante $\arg\max$
    \item No puede aprender una \textbf{política estocástica} (Stochastic Policy), y en algunos escenarios la aleatoriedad es beneficiosa
\end{itemize}

\begin{warningbox}{¿Cuándo no es aplicable Q-Learning?}
Cuando el espacio de acción es continuo (como los ángulos de las articulaciones de un robot, o la fuerza del acelerador y el freno de un automóvil), Q-Learning no puede aplicarse directamente, porque necesita calcular el valor Q para \textbf{todas las acciones posibles} y tomar el máximo. En un espacio continuo hay infinitas acciones, lo cual es computacionalmente inviable. En ese caso hay que recurrir a los métodos de Policy Gradient.
\end{warningbox}

\subsection{Resumen de la sección}

DQN aproxima la Q-function mediante una red neuronal profunda, con lo que logra un aprendizaje de extremo a extremo desde los píxeles sin procesar hasta la decisión. Su entrenamiento se basa en construir el target con la ecuación de Bellman y actualizar los parámetros de la red minimizando la Q-Loss. DQN alcanzó resultados que superan el nivel humano en juegos de Atari, pero, al estar limitado a espacios de acción discretos y políticas deterministas, no puede manejar directamente problemas más complejos como el control continuo.

